% rel_pnofs.tex
%
% Explicit reconstructed-2-RDM elements [cf. Eqs. (98a)-(98f) of
% SciPost Chem. 1, 004 (2022)] for each of the relativistic PNOF5,
% PNOF7, PNOF7s and GNOF functionals, together with the complete,
% step-by-step algebraic derivation showing that substitution of
% these elements into Eq. (64) of the source paper reproduces the
% published energy expressions of Sec. 3.5.3 [Eqs. (72)-(76)].

\documentclass[11pt,a4paper]{article}

\usepackage{amsmath,amssymb,mathtools,bm}
\usepackage{geometry}
\usepackage{booktabs}
\usepackage{hyperref}
\usepackage{microtype}
\usepackage{enumitem}

\geometry{margin=1in}
\hypersetup{colorlinks=true,linkcolor=blue,citecolor=blue,urlcolor=blue}

\newcommand{\Dtwo}{{}^2D}
\newcommand{\PiMat}{\Pi}
\newcommand{\DeltaMat}{\Delta}
\newcommand{\Om}{\Omega}
\newcommand{\hi}{h_i}
\newcommand{\hj}{h_j}

\title{Explicit Reconstructed 2-RDM Elements and Full Derivation of the\\
Relativistic PNOF5, PNOF7, PNOF7s, and GNOF Energy Functionals}
\author{}
\date{}

\begin{document}
\maketitle

\begin{abstract}
Starting from the general no-pair relativistic 2-RDM reconstruction of
Appendix D of SciPost Chem.\ \textbf{1}, 004 (2022) [Eqs.~(97a)--(98f)],
we write down, fully explicitly and separately for the relativistic
PNOF5, PNOF7, PNOF7s, and GNOF functionals, every 2-RDM element
$\Dtwo^{kl}_{ij}$ entering the natural-orbital energy expression
Eq.~(64) of that paper. We then substitute these explicit elements
into Eq.~(64), evaluate every one of its six blocks in full detail,
and show that the result is identical, term by term, to the published
energy expressions Eqs.~(72)--(76) of Sec.~3.5.3. The no-pair
Dirac--Hartree--Fock limit is recovered as a corollary.
\end{abstract}

\tableofcontents

\section{Scope, conventions, and notation}
\label{sec:scope}

We work throughout within the no-pair approximation and impose
Kramers' pairing symmetry, restricting attention to `spin-compensated'
systems, exactly as in Sec.~3.5 of the source paper. Natural spinors
are grouped into Kramers pairs $(i,\bar i)$, with $i$ running over
`half of the positive-energy spinors' and $\bar i$ its time-reversed
partner. Spin compensation means
\begin{equation}
 n_{\bar i}=n_i \qquad\text{for all } i.
 \label{eq:kramers-occ}
\end{equation}
The spinor space is partitioned into mutually disjoint subspaces
(geminals),
\begin{equation}
 \Om=\bigcup_p \Om_p, \qquad \Om_p\cap\Om_q=\varnothing \ (p\neq q),
 \qquad \sum_p n_p =\tfrac12 N_e ,
 \label{eq:partition}
\end{equation}
with $n_p\equiv n_i$ for the reference (`principal') orbital of
$\Om_p$ and $h_p=1-n_p$.

We use the same integral symbols as the source paper,
\begin{align}
 h_{ii}&=\int d\mathbf r\,d\mathbf r'\,\widetilde\chi_i^\dagger(\mathbf r')
 \big[\delta(\mathbf r-\mathbf r')\widehat T_D(\mathbf r)+v_{\rm ext}^{\rm nl}(\mathbf r,\mathbf r')\big]\widetilde\chi_i(\mathbf r),\\
 J_{ij}&=\langle ij|ij\rangle, &
 J^G_{ij}&=\langle ij|\bm\alpha\!\cdot\!\bm\alpha|ij\rangle,\\
 K_{ij}&=\langle ij|ji\rangle, &
 K^G_{ij}&=\langle ij|\bm\alpha\!\cdot\!\bm\alpha|ji\rangle,\\
 L_{ij}&=\langle \bar i j|j\bar i\rangle, &
 L^G_{ij}&=\langle \bar i j|\bm\alpha\!\cdot\!\bm\alpha|j\bar i\rangle,
\end{align}
with the paper's stated property $L_{ii}=0$ (while $L^G_{ii}\neq0$ in
general).

The cumulant reconstruction of Appendix D uses two auxiliary real,
symmetric matrices, $\DeltaMat_{ij}$ and $\PiMat_{ij}\equiv
\PiMat_{i\bar i,j\bar j}$ (the four possible bar-orderings of the
latter are all numerically equal, as shown explicitly below
Eq.~(97e) of the source paper). Following the prescription used to
build Eqs.~(72)--(76), we adopt
\begin{equation}
 \DeltaMat_{ij}=\DeltaMat_{i\bar j}=\DeltaMat_{\bar ij}=\DeltaMat_{\bar i\bar j},
 \qquad
 \DeltaMat_{ii}=n_i^2,
 \qquad
 \PiMat_{ii}=n_i,
 \label{eq:delta-pi-diag}
\end{equation}
and, for the off-diagonal part of $\DeltaMat$, the subspace ansatz
\begin{equation}
 \DeltaMat_{ij}=
 \begin{cases}
 n_i n_j, & i,j\in\Om_p \ (i\neq j),\\[1mm]
 0, & i\in\Om_q,\ j\in\Om_p,\quad p\neq q.
 \end{cases}
 \label{eq:delta-subspaces}
\end{equation}
This ansatz is the \emph{same for all four functionals}; only
$\PiMat^{\rm inter}$ differs between PNOF5, PNOF7, PNOF7s, and GNOF.
For compactness define
\begin{equation}
 A_{ij}\equiv\frac{n_in_j-\DeltaMat_{ij}}{2}
 \;\Longrightarrow\;
 A_{ij}=
 \begin{cases}
 0, & i,j\in\Om_p,\\[1mm]
 \dfrac{n_in_j}{2}, & i\in\Om_q,\ j\in\Om_p\ (p\neq q).
 \end{cases}
 \label{eq:Aij}
\end{equation}

The intra-subspace part of $\PiMat$ is likewise common to all four
functionals [Eq.~(74) of the source paper]:
\begin{equation}
 \PiMat^{\rm intra}_{ij}=
 \begin{cases}
 -\sqrt{n_in_j}, & i=p\ \text{or}\ j=p \ (\text{principal orbital of }\Om_p),\\[1mm]
 +\sqrt{n_in_j}, & \text{otherwise},
 \end{cases}
 \qquad i,j\in\Om_p,\ i\neq j.
 \label{eq:Pi-intra}
\end{equation}
The inter-subspace part, $\PiMat^{\rm inter}_{ij;pq}$ for
$i\in\Om_q,\,j\in\Om_p,\,p\neq q$, is what distinguishes the four
functionals (Table 3 of the source paper) and is given explicitly in
Sec.~\ref{sec:functionals} below.

\section{The general reconstructed 2-RDM [Eqs.~(98a)--(98f)]}
\label{sec:general98}

For reference, the general (functional-independent) reconstruction of
Appendix D reads
\begin{align}
\Dtwo^{kl}_{ij} &=
 \frac{n_in_j-\DeltaMat_{ij}}{2}\delta_{ik}\delta_{jl}
 -\frac{n_in_j-\DeltaMat_{ij}}{2}\delta_{il}\delta_{jk},
 \label{eq:98a}\\
\Dtwo^{\bar k\bar l}_{\bar i\bar j} &=
 \frac{n_{\bar i}n_{\bar j}-\DeltaMat_{\bar i\bar j}}{2}\delta_{ik}\delta_{jl}
 -\frac{n_{\bar i}n_{\bar j}-\DeltaMat_{\bar i\bar j}}{2}\delta_{il}\delta_{jk},
 \label{eq:98b}\\
\Dtwo^{k\bar l}_{i\bar j} &=
 \frac{n_in_{\bar j}-\DeltaMat_{i\bar j}}{2}\delta_{ik}\delta_{jl}
 +\frac{\PiMat_{i\bar i,k\bar k}}{2}\delta_{ij}\delta_{kl},
 \label{eq:98c}\\
\Dtwo^{\bar lk}_{\bar ji} &=
 \frac{n_in_{\bar j}-\DeltaMat_{\bar ji}}{2}\delta_{ik}\delta_{jl}
 +\frac{\PiMat_{\bar ii,\bar k\bar k}}{2}\delta_{ij}\delta_{kl},
 \label{eq:98d}\\
\Dtwo^{k\bar l}_{\bar ji} &=
 -\frac{n_in_{\bar j}-\DeltaMat_{\bar ji}}{2}\delta_{ik}\delta_{jl}
 -\frac{\PiMat_{\bar ii,k\bar k}}{2}\delta_{ij}\delta_{kl},
 \label{eq:98e}\\
\Dtwo^{\bar lk}_{i\bar j} &=
 -\frac{n_in_{\bar j}-\DeltaMat_{i\bar j}}{2}\delta_{ik}\delta_{jl}
 -\frac{\PiMat_{i\bar i,\bar k\bar k}}{2}\delta_{ij}\delta_{kl}.
 \label{eq:98f}
\end{align}
Using spin compensation \eqref{eq:kramers-occ}, the $\DeltaMat$
symmetry \eqref{eq:delta-pi-diag}, and the equality of all
bar-orderings of $\PiMat$, every element that enters Eq.~(64) reduces
to one of only two closed forms, valid for \emph{any} $i,k$ (equal or
not):
\begin{align}
\Dtwo^{k\bar k}_{i\bar i} &= \frac{n_i^2-\DeltaMat_{ii}}{2}\delta_{ik}
 +\frac{\PiMat_{ik}}{2}
 \;=\;\frac{\PiMat_{ik}}{2},
 \label{eq:unified-plus}\\
\Dtwo^{i\bar i}_{\bar kk} &= -\frac{n_k^2-\DeltaMat_{kk}}{2}\delta_{ik}
 -\frac{\PiMat_{ik}}{2}
 \;=\;-\frac{\PiMat_{ik}}{2},
 \label{eq:unified-minus}
\end{align}
where the diagonal-$\DeltaMat$ convention $\DeltaMat_{ii}=n_i^2$ kills
the first term identically (this is what makes the diagonal case,
$i=k$, consistent with the off-diagonal one via
$\PiMat_{ii}=n_i$). Equations \eqref{eq:unified-plus}--\eqref{eq:unified-minus}
are the compact `Eq.~98' library used repeatedly below.

The remaining `Coulomb/exchange-type' elements needed are, for
$i\neq j$,
\begin{equation}
\Dtwo^{ij}_{ij}=\Dtwo^{\bar ij}_{\bar ij}=\Dtwo^{i\bar j}_{i\bar j}
=\Dtwo^{\bar i\bar j}_{\bar i\bar j}=A_{ij},
\qquad
\Dtwo^{ji}_{ij}=\Dtwo^{\bar i\bar j}_{\bar j\bar i}=-A_{ij},
\label{eq:coulomb-elems}
\end{equation}
and the four `spin-flip' ($L$-type) elements, for $i\neq j$,
\begin{equation}
\Dtwo^{j\bar i}_{\bar ij}=\Dtwo^{i\bar j}_{\bar ji}
=\Dtwo^{\bar ij}_{j\bar i}=\Dtwo^{\bar ji}_{i\bar j}=-A_{ij}.
\label{eq:lflip-elems}
\end{equation}
For $i=j$ the pure-unbarred/pure-barred elements
$\Dtwo^{ii}_{ii}=\Dtwo^{\bar i\bar i}_{\bar i\bar i}=0$ vanish
identically by antisymmetry, while $\Dtwo^{i\bar i}_{i\bar i}=\Dtwo^{\bar ii}_{\bar ii}=n_i/2$
follow from \eqref{eq:unified-plus} at $i=k$, and
$\Dtwo^{i\bar i}_{\bar ii}=\Dtwo^{\bar ii}_{i\bar i}=-n_i/2$ follow
from \eqref{eq:unified-minus} at $i=k$.

\section{Explicit 2-RDM elements for each functional}
\label{sec:functionals}

Because \eqref{eq:Aij}, \eqref{eq:coulomb-elems}, and
\eqref{eq:lflip-elems} depend on $\DeltaMat$ only, they are
\emph{identical for all four functionals} once the intra/inter split
\eqref{eq:delta-subspaces} is applied. The functional identity enters
only through $\PiMat_{ik}$ in \eqref{eq:unified-plus}--\eqref{eq:unified-minus},
via $\PiMat^{\rm intra}$ [Eq.~\eqref{eq:Pi-intra}, common to all] and
$\PiMat^{\rm inter}$ [functional-dependent]. We now list, functional
by functional, the complete set of explicit 2-RDM elements.

\subsection*{Common intra-subspace elements ($i,j\in\Om_p$, $i\neq j$)}
\begin{align}
\Dtwo^{ij}_{ij}=\Dtwo^{\bar ij}_{\bar ij}=\Dtwo^{i\bar j}_{i\bar j}
=\Dtwo^{\bar i\bar j}_{\bar i\bar j}&=0, &
\Dtwo^{ji}_{ij}=\Dtwo^{\bar i\bar j}_{\bar j\bar i}&=0,
\label{eq:intra-coulomb}\\
\Dtwo^{j\bar i}_{\bar ij}=\Dtwo^{i\bar j}_{\bar ji}
=\Dtwo^{\bar ij}_{j\bar i}=\Dtwo^{\bar ji}_{i\bar j}&=0,
\label{eq:intra-lflip}\\
\Dtwo^{j\bar j}_{i\bar i}=\Dtwo^{\bar ii}_{\bar jj}
&=\frac{\PiMat^{\rm intra}_{ij}}{2}
=\mp\frac{\sqrt{n_in_j}}{2},
\label{eq:intra-piK}\\
\Dtwo^{i\bar i}_{\bar jj}=\Dtwo^{j\bar j}_{\bar ii}
=\Dtwo^{\bar ii}_{j\bar j}=\Dtwo^{\bar jj}_{i\bar i}
&=-\frac{\PiMat^{\rm intra}_{ij}}{2}
=\pm\frac{\sqrt{n_in_j}}{2},
\label{eq:intra-piL}
\end{align}
(upper sign if $i$ or $j$ is the principal orbital of $\Om_p$, lower
sign otherwise), all vanishing whenever $\Delta_{ij}=n_in_j$ makes
$A_{ij}=0$, since intra-subspace pairs never contribute a
Coulomb/exchange-type or $L$-flip-type 2-RDM element -- only the
$\PiMat^{\rm intra}$ pairing term survives.

\subsection*{Common diagonal elements ($i=j$, any functional)}
\begin{equation}
\Dtwo^{ii}_{ii}=\Dtwo^{\bar i\bar i}_{\bar i\bar i}=0,\qquad
\Dtwo^{i\bar i}_{i\bar i}=\Dtwo^{\bar ii}_{\bar ii}=\frac{n_i}{2},\qquad
\Dtwo^{i\bar i}_{\bar ii}=\Dtwo^{\bar ii}_{i\bar i}=-\frac{n_i}{2}.
\label{eq:diag-common}
\end{equation}

\subsection*{Inter-subspace elements ($i\in\Om_q$, $j\in\Om_p$, $q\neq p$)}
Here $A_{ij}=n_in_j/2$, so
\begin{align}
\Dtwo^{ij}_{ij}=\Dtwo^{\bar ij}_{\bar ij}=\Dtwo^{i\bar j}_{i\bar j}
=\Dtwo^{\bar i\bar j}_{\bar i\bar j}&=\frac{n_in_j}{2}, &
\Dtwo^{ji}_{ij}=\Dtwo^{\bar i\bar j}_{\bar j\bar i}&=-\frac{n_in_j}{2},
\label{eq:inter-coulomb}\\
\Dtwo^{j\bar i}_{\bar ij}=\Dtwo^{i\bar j}_{\bar ji}
=\Dtwo^{\bar ij}_{j\bar i}=\Dtwo^{\bar ji}_{i\bar j}&=-\frac{n_in_j}{2},
\label{eq:inter-lflip}\\
\Dtwo^{j\bar j}_{i\bar i}=\Dtwo^{\bar ii}_{\bar jj}
&=\frac{\PiMat^{\rm inter}_{ij;pq}}{2},
\label{eq:inter-piK}\\
\Dtwo^{i\bar i}_{\bar jj}=\Dtwo^{j\bar j}_{\bar ii}
=\Dtwo^{\bar ii}_{j\bar j}=\Dtwo^{\bar jj}_{i\bar i}
&=-\frac{\PiMat^{\rm inter}_{ij;pq}}{2}.
\label{eq:inter-piL}
\end{align}
These are common in \emph{form} to all four functionals; only the
value of $\PiMat^{\rm inter}_{ij;pq}$ in
\eqref{eq:inter-piK}--\eqref{eq:inter-piL} changes. We now give this
value, and hence the fully explicit inter-subspace 2-RDM elements,
functional by functional.

\subsubsection{PNOF5}
\begin{equation}
\PiMat^{\rm inter,\,PNOF5}_{ij;pq}=0
\;\;\Longrightarrow\;\;
\Dtwo^{j\bar j}_{i\bar i}=\Dtwo^{\bar ii}_{\bar jj}
=\Dtwo^{i\bar i}_{\bar jj}=\Dtwo^{j\bar j}_{\bar ii}
=\Dtwo^{\bar ii}_{j\bar j}=\Dtwo^{\bar jj}_{i\bar i}=0.
\label{eq:98-pnof5}
\end{equation}
For PNOF5 the inter-subspace $2$-RDM therefore contains \emph{only}
the ordinary Coulomb/exchange/$L$-flip elements
\eqref{eq:inter-coulomb}--\eqref{eq:inter-lflip}; there is no
inter-subspace pairing contribution.

\subsubsection{PNOF7}
\begin{equation}
\PiMat^{\rm inter,\,PNOF7}_{ij;pq}=-\sqrt{n_in_j\hi\hj}
\;\;\Longrightarrow\;\;
\begin{aligned}
\Dtwo^{j\bar j}_{i\bar i}&=\Dtwo^{\bar ii}_{\bar jj}
=-\tfrac12\sqrt{n_in_j\hi\hj},\\
\Dtwo^{i\bar i}_{\bar jj}&=\Dtwo^{j\bar j}_{\bar ii}
=\Dtwo^{\bar ii}_{j\bar j}=\Dtwo^{\bar jj}_{i\bar i}
=+\tfrac12\sqrt{n_in_j\hi\hj}.
\end{aligned}
\label{eq:98-pnof7}
\end{equation}

\subsubsection{PNOF7s}
\begin{equation}
\PiMat^{\rm inter,\,PNOF7s}_{ij;pq}=-4n_in_j\hi\hj
\;\;\Longrightarrow\;\;
\begin{aligned}
\Dtwo^{j\bar j}_{i\bar i}&=\Dtwo^{\bar ii}_{\bar jj}
=-2n_in_j\hi\hj,\\
\Dtwo^{i\bar i}_{\bar jj}&=\Dtwo^{j\bar j}_{\bar ii}
=\Dtwo^{\bar ii}_{j\bar j}=\Dtwo^{\bar jj}_{i\bar i}
=+2n_in_j\hi\hj.
\end{aligned}
\label{eq:98-pnof7s}
\end{equation}

\subsubsection{GNOF}
Define, with $h_p=1-n_p$ (the hole of the principal orbital of
$\Om_p$),
\begin{equation}
 n_i^d=n_i\,\frac{h_p^d}{h_p},
 \qquad
 h_p^d=h_p\exp\!\left[-\Big(\frac{h_p}{0.02\sqrt2}\Big)^2\right].
\end{equation}
Then, for $i\in\Om_q$, $j\in\Om_p$, $q\neq p$,
\begin{equation}
\PiMat^{\rm inter,\,GNOF}_{ij;pq}=
\begin{cases}
n_i^dn_j^d-\sqrt{n_in_j\hi\hj}-\sqrt{n_i^dn_j^d}, & i>q,\ j=p\ \text{or}\ i=q,\ j>p,\\[1mm]
n_i^dn_j^d-\sqrt{n_in_j\hi\hj}+\sqrt{n_i^dn_j^d}, & i>q,\ j>p,\\[1mm]
0, & i=q,\ j=p,
\end{cases}
\label{eq:Pi-gnof}
\end{equation}
so that the fully explicit GNOF inter-subspace elements are
\begin{equation}
\Dtwo^{j\bar j}_{i\bar i}=\Dtwo^{\bar ii}_{\bar jj}
=\frac{\PiMat^{\rm inter,\,GNOF}_{ij;pq}}{2},
\qquad
\Dtwo^{i\bar i}_{\bar jj}=\Dtwo^{j\bar j}_{\bar ii}
=\Dtwo^{\bar ii}_{j\bar j}=\Dtwo^{\bar jj}_{i\bar i}
=-\frac{\PiMat^{\rm inter,\,GNOF}_{ij;pq}}{2},
\label{eq:98-gnof}
\end{equation}
with $\PiMat^{\rm inter,\,GNOF}_{ij;pq}$ given piecewise by
\eqref{eq:Pi-gnof}.

Equations \eqref{eq:98-pnof5}, \eqref{eq:98-pnof7},
\eqref{eq:98-pnof7s}, and \eqref{eq:98-gnof} -- together with the
common elements \eqref{eq:intra-coulomb}--\eqref{eq:diag-common} and
\eqref{eq:inter-coulomb}--\eqref{eq:inter-lflip} -- constitute the
complete, explicit reconstructed 2-RDM (`Eq.~98') for each of the
four functionals. We now show, step by step, that inserting these
elements into Eq.~(64) reproduces Eqs.~(72)--(76).

\section{The energy expression, Eq.~(64)}
\label{sec:E64}

The np-ReNOFT energy under Kramers symmetry (Eq.~(64) of the source
paper) is
\begin{align}
E={}&\sum_i h_{ii}(n_i+n_{\bar i})
\nonumber\\
&+\sum_{i,j}\Big(\Dtwo^{ij}_{ij}+\Dtwo^{\bar ij}_{\bar ij}
+\Dtwo^{i\bar j}_{i\bar j}+\Dtwo^{\bar i\bar j}_{\bar i\bar j}\Big)J_{ij}
\nonumber\\
&-\sum_{i,j}\Big(\Dtwo^{ij}_{ij}-\Dtwo^{\bar ij}_{\bar ij}
-\Dtwo^{i\bar j}_{i\bar j}+\Dtwo^{\bar i\bar j}_{\bar i\bar j}\Big)J^G_{ij}
\nonumber\\
&+\sum_{i,j}\Big[\Dtwo^{ji}_{ij}+\Dtwo^{\bar i\bar j}_{\bar j\bar i}\Big](K_{ij}-K^G_{ij})
\nonumber\\
&+\frac12\sum_{i,j}\Big[\Dtwo^{j\bar i}_{\bar ij}+\Dtwo^{i\bar j}_{\bar ji}
+\Dtwo^{\bar ij}_{j\bar i}+\Dtwo^{\bar ji}_{i\bar j}\Big](L_{ij}-L^G_{ij})
\nonumber\\
&+\sum_{i\neq j}\Big[\Dtwo^{j\bar j}_{i\bar i}+\Dtwo^{\bar ii}_{\bar jj}\Big](K_{ij}+K^G_{ij})
\nonumber\\
&-\frac12\sum_{i\neq j}\Big[\Dtwo^{i\bar i}_{\bar jj}+\Dtwo^{j\bar j}_{\bar ii}
+\Dtwo^{\bar ii}_{j\bar j}+\Dtwo^{\bar jj}_{i\bar i}\Big](L_{ij}+L^G_{ij}).
\label{eq:E64}
\end{align}
The six sums (call them, in order, the $h$-, $J$-, $J^G$-, $K$-,
$L$-, $\PiMat K$-, and $\PiMat L$-blocks) are evaluated one at a time
below, generically in $A_{ij}$ and $\PiMat_{ij}$; the functional
identity is inserted only at the very end.

\subsection{The $h$-block}
By spin compensation \eqref{eq:kramers-occ},
\begin{equation}
\sum_i h_{ii}(n_i+n_{\bar i})=2\sum_i n_ih_{ii}.
\label{eq:hblock}
\end{equation}

\subsection{The Coulomb ($J$) block}
Define $C^J_{ij}=\Dtwo^{ij}_{ij}+\Dtwo^{\bar ij}_{\bar ij}+\Dtwo^{i\bar j}_{i\bar j}+\Dtwo^{\bar i\bar j}_{\bar i\bar j}$.
For $i\neq j$, using \eqref{eq:coulomb-elems},
\begin{equation}
C^J_{ij}=4A_{ij}=2(n_in_j-\DeltaMat_{ij}).
\label{eq:CJ-offdiag}
\end{equation}
For $i=j$, using \eqref{eq:diag-common},
\begin{equation}
C^J_{ii}=0+0+\frac{n_i}{2}+\frac{n_i}{2}=n_i.
\label{eq:CJ-diag}
\end{equation}
Hence
\begin{equation}
E_J=\sum_i n_iJ_{ii}+\sum_{i\neq j}2(n_in_j-\DeltaMat_{ij})J_{ij}.
\label{eq:EJ}
\end{equation}

\subsection{The Gaunt ($J^G$) block}
Define $C^{J^G}_{ij}=\Dtwo^{ij}_{ij}-\Dtwo^{\bar ij}_{\bar ij}-\Dtwo^{i\bar j}_{i\bar j}+\Dtwo^{\bar i\bar j}_{\bar i\bar j}$.
For $i\neq j$, all four elements equal $A_{ij}$, so
\begin{equation}
C^{J^G}_{ij}=A_{ij}-A_{ij}-A_{ij}+A_{ij}=0.
\end{equation}
For $i=j$,
\begin{equation}
C^{J^G}_{ii}=0-\frac{n_i}{2}-\frac{n_i}{2}+0=-n_i .
\end{equation}
Since Eq.~\eqref{eq:E64} contains $-C^{J^G}_{ij}J^G_{ij}$,
\begin{equation}
E_{J^G}=\sum_i n_iJ^G_{ii}.
\label{eq:EJG}
\end{equation}
No inter- or intra-subspace ($i\neq j$) Gaunt-Coulomb term survives:
this is why Eqs.~(73) and (75) contain $J^G_{ii}$ (diagonal only) and
never $J^G_{ij}$ ($i\ne j$).

\subsection{The exchange ($K$) block}
Define $C^K_{ij}=\Dtwo^{ji}_{ij}+\Dtwo^{\bar i\bar j}_{\bar j\bar i}$.
For $i\neq j$, using \eqref{eq:coulomb-elems},
\begin{equation}
C^K_{ij}=-2A_{ij}=-(n_in_j-\DeltaMat_{ij}).
\label{eq:CK}
\end{equation}
For $i=j$ both terms vanish by antisymmetry. Hence
\begin{equation}
E_K=-\sum_{i\neq j}(n_in_j-\DeltaMat_{ij})(K_{ij}-K^G_{ij}).
\label{eq:EK}
\end{equation}

\subsection{The spin-flip ($L$) block}
Define $C^L_{ij}=\Dtwo^{j\bar i}_{\bar ij}+\Dtwo^{i\bar j}_{\bar ji}+\Dtwo^{\bar ij}_{j\bar i}+\Dtwo^{\bar ji}_{i\bar j}$
(note that in Eq.~\eqref{eq:E64} this sum runs over \emph{all} $i,j$,
including $i=j$, unlike the $K$- and $\PiMat$-blocks).

For $i\neq j$, using \eqref{eq:lflip-elems},
\begin{equation}
C^L_{ij}=-4A_{ij}=-2(n_in_j-\DeltaMat_{ij}) .
\end{equation}
With the $\tfrac12$ prefactor of Eq.~\eqref{eq:E64},
\begin{equation}
E_L^{i\neq j}=-\sum_{i\neq j}(n_in_j-\DeltaMat_{ij})(L_{ij}-L^G_{ij}).
\label{eq:EL-offdiag}
\end{equation}

For $i=j$, using \eqref{eq:diag-common}, all four terms collapse
pairwise into $2\Dtwo^{i\bar i}_{\bar ii}+2\Dtwo^{\bar ii}_{i\bar i}
=2(-n_i/2)+2(-n_i/2)=-2n_i$, so
\begin{equation}
C^L_{ii}=-2n_i .
\label{eq:CLdiag}
\end{equation}
With the $\tfrac12$ prefactor and $L_{ii}=0$ (source paper),
\begin{equation}
\frac12C^L_{ii}(L_{ii}-L^G_{ii})
=-n_i(0-L^G_{ii})
=+n_iL^G_{ii}.
\label{eq:ELdiag}
\end{equation}
This diagonal contribution, $\sum_in_iL^G_{ii}$, is precisely the
term that appears together with $2h_{ii}+J_{ii}+J^G_{ii}$ in the
one-body-like sum of Eqs.~(72)--(76), even though $L_{ii}=0$
identically.

\subsection{The $\PiMat(K+K^G)$ block}
For $i\neq j$, using \eqref{eq:unified-plus},
\begin{equation}
\Dtwo^{j\bar j}_{i\bar i}=\Dtwo^{\bar ii}_{\bar jj}=\frac{\PiMat_{ij}}{2}
\;\Longrightarrow\;
E_{\PiMat K}=\sum_{i\neq j}\PiMat_{ij}(K_{ij}+K^G_{ij}).
\label{eq:EpiK}
\end{equation}

\subsection{The $\PiMat(L+L^G)$ block}
For $i\neq j$, using \eqref{eq:unified-minus},
\begin{equation}
\Dtwo^{i\bar i}_{\bar jj}=\Dtwo^{j\bar j}_{\bar ii}=\Dtwo^{\bar ii}_{j\bar j}
=\Dtwo^{\bar jj}_{i\bar i}=-\frac{\PiMat_{ij}}{2}
\;\Longrightarrow\;
-\frac12\big(-4\cdot\tfrac{\PiMat_{ij}}{2}\big)=\PiMat_{ij},
\end{equation}
so
\begin{equation}
E_{\PiMat L}=\sum_{i\neq j}\PiMat_{ij}(L_{ij}+L^G_{ij}).
\label{eq:EpiL}
\end{equation}
Combining \eqref{eq:EpiK} and \eqref{eq:EpiL},
\begin{equation}
E_\PiMat=\sum_{i\neq j}\PiMat_{ij}\big(K_{ij}+K^G_{ij}+L_{ij}+L^G_{ij}\big).
\label{eq:Epi}
\end{equation}

\section{The master formula}
\label{sec:master}

Collecting \eqref{eq:hblock}, \eqref{eq:EJ}, \eqref{eq:EJG},
\eqref{eq:EK}, \eqref{eq:EL-offdiag}, \eqref{eq:ELdiag}, and
\eqref{eq:Epi}, and grouping diagonal versus off-diagonal terms,
\begin{equation}
\boxed{
\begin{aligned}
E={}&\sum_in_i\big(2h_{ii}+J_{ii}+J^G_{ii}+L^G_{ii}\big)\\
&+\sum_{i\neq j}(n_in_j-\DeltaMat_{ij})\Big[2J_{ij}-(K_{ij}-K^G_{ij})-(L_{ij}-L^G_{ij})\Big]\\
&+\sum_{i\neq j}\PiMat_{ij}\Big[K_{ij}+K^G_{ij}+L_{ij}+L^G_{ij}\Big].
\end{aligned}
}
\label{eq:master}
\end{equation}
This identity holds for \emph{any} valid choice of $\DeltaMat$ and
$\PiMat$ satisfying the conventions of Sec.~\ref{sec:scope}; it is
independent of which of the four Piris-type functionals is used.

\section{Separation into intra- and inter-subspace terms}
\label{sec:separation}

Insert the subspace ansatz \eqref{eq:delta-subspaces} into
\eqref{eq:master}.

\subsection*{Intra-subspace ($i,j\in\Om_p$)}
Here $n_in_j-\DeltaMat_{ij}=0$, so the second sum in
\eqref{eq:master} drops out entirely and only $\PiMat^{\rm intra}$
survives:
\begin{equation}
\boxed{
E_p=\sum_{i\in\Om_p}n_i\big(2h_{ii}+J_{ii}+J^G_{ii}+L^G_{ii}\big)
+\sum_{\substack{i,j\in\Om_p\\i\neq j}}\PiMat^{\rm intra}_{ij}\big(K_{ij}+K^G_{ij}+L_{ij}+L^G_{ij}\big).
}
\label{eq:Ep}
\end{equation}
This is \emph{identical} to Eq.~(73) of the source paper.

\subsection*{Inter-subspace ($i\in\Om_q$, $j\in\Om_p$, $q\neq p$)}
Here $n_in_j-\DeltaMat_{ij}=n_in_j$, so
\begin{equation}
\boxed{
E_{qp}=\sum_{i\in\Om_q}\sum_{j\in\Om_p}n_in_j\Big[2J_{ij}-(K_{ij}-K^G_{ij})-(L_{ij}-L^G_{ij})\Big]
+\sum_{i\in\Om_q}\sum_{j\in\Om_p}\PiMat^{\rm inter}_{ij;pq}\big(K_{ij}+K^G_{ij}+L_{ij}+L^G_{ij}\big).
}
\label{eq:Eqp}
\end{equation}
This is \emph{identical} to Eq.~(75) of the source paper. Since
$E=\sum_pE_p+\sum_{q\neq p}E_{qp}$, Eq.~(72) is recovered exactly.

Equations \eqref{eq:Ep}--\eqref{eq:Eqp} are still generic in
$\PiMat^{\rm inter}_{ij;pq}$: they hold for \emph{any} of the four
functionals. We now specialize.

\section{Recovery of each functional's final energy expression}
\label{sec:recovery}

\subsection{PNOF5}
Substituting $\PiMat^{\rm inter}_{ij;pq}=0$ [Eq.~\eqref{eq:98-pnof5}]
into \eqref{eq:Eqp} kills the second sum identically:
\begin{equation}
\boxed{
E_{qp}^{\rm PNOF5}=\sum_{i\in\Om_q}\sum_{j\in\Om_p}n_in_j\Big[2J_{ij}-(K_{ij}-K^G_{ij})-(L_{ij}-L^G_{ij})\Big].
}
\label{eq:PNOF5-final}
\end{equation}
Together with the (unchanged) intra-subspace term \eqref{eq:Ep} using
$\PiMat^{\rm intra}$ of \eqref{eq:Pi-intra}, this reproduces the
relativistic PNOF5 functional exactly.

\subsection{PNOF7}
Substituting $\PiMat^{\rm inter}_{ij;pq}=-\sqrt{n_in_j\hi\hj}$
[Eq.~\eqref{eq:98-pnof7}] into \eqref{eq:Eqp},
\begin{equation}
\boxed{
E_{qp}^{\rm PNOF7}=\sum_{i\in\Om_q}\sum_{j\in\Om_p}n_in_j\Big[2J_{ij}-(K_{ij}-K^G_{ij})-(L_{ij}-L^G_{ij})\Big]
-\sum_{i\in\Om_q}\sum_{j\in\Om_p}\sqrt{n_in_j\hi\hj}\big(K_{ij}+K^G_{ij}+L_{ij}+L^G_{ij}\big).
}
\label{eq:PNOF7-final}
\end{equation}
This is exactly the relativistic PNOF7 term of Eq.~(75).

\subsection{PNOF7s}
Substituting $\PiMat^{\rm inter}_{ij;pq}=-4n_in_j\hi\hj$
[Eq.~\eqref{eq:98-pnof7s}] into \eqref{eq:Eqp},
\begin{equation}
\boxed{
E_{qp}^{\rm PNOF7s}=\sum_{i\in\Om_q}\sum_{j\in\Om_p}n_in_j\Big[2J_{ij}-(K_{ij}-K^G_{ij})-(L_{ij}-L^G_{ij})\Big]
-4\sum_{i\in\Om_q}\sum_{j\in\Om_p}n_in_j\hi\hj\big(K_{ij}+K^G_{ij}+L_{ij}+L^G_{ij}\big).
}
\label{eq:PNOF7s-final}
\end{equation}

\subsection{GNOF}
Substituting the piecewise $\PiMat^{\rm inter,\,GNOF}_{ij;pq}$ of
\eqref{eq:Pi-gnof} into \eqref{eq:Eqp} gives, with no further change
to the reconstructed-2-RDM algebra,
\begin{equation}
\boxed{
E_{qp}^{\rm GNOF}=\sum_{i\in\Om_q}\sum_{j\in\Om_p}n_in_j\Big[2J_{ij}-(K_{ij}-K^G_{ij})-(L_{ij}-L^G_{ij})\Big]
+\sum_{i\in\Om_q}\sum_{j\in\Om_p}\PiMat^{\rm inter,\,GNOF}_{ij;pq}\big(K_{ij}+K^G_{ij}+L_{ij}+L^G_{ij}\big),
}
\label{eq:GNOF-final}
\end{equation}
i.e.\ exactly the relativistic GNOF term of Eq.~(75)--(76).

In every case,
\begin{equation}
E^{x}=\sum_pE_p+\sum_{q\neq p}E_{qp}^{x},\qquad x\in\{{\rm PNOF5,\,PNOF7,\,PNOF7s,\,GNOF}\},
\end{equation}
with $E_p$ common [Eq.~\eqref{eq:Ep}] and $E_{qp}^{x}$ given by
\eqref{eq:PNOF5-final}--\eqref{eq:GNOF-final}: these are precisely
Eqs.~(72) and (75) of the source paper, specialized functional by
functional.

\section{The no-pair Dirac--Hartree--Fock limit}
\label{sec:HFlimit}

As a corollary, let $n_i\in\{0,1\}$ for every $i$ (idempotent
1-RDM). Then $h_i=1-n_i\in\{0,1\}$ as well, and for \emph{any} pair
$i,j$ (same or different subspace) either $n_i=0$, $n_j=0$, $h_i=0$,
or $h_j=0$. Hence
\begin{equation}
\sqrt{n_in_j}=0\ \text{or}\ 1\ \text{trivially},\qquad
\sqrt{n_in_j\hi\hj}=0,\qquad n_in_j\hi\hj=0,\qquad
n_i^dn_j^d-\sqrt{n_in_j\hi\hj}\pm\sqrt{n_i^dn_j^d}\to0,
\end{equation}
because in the strict single-determinant limit each subspace $\Om_p$
contains exactly one occupied orbital ($n_i=1$) and the rest
unoccupied ($n_j=0$); this kills $\PiMat^{\rm intra}_{ij}$ as well
(one of $n_i,n_j$ vanishes for any intra-subspace pair with $i\ne j$
in this limit). Consequently every $\PiMat$-dependent term in
\eqref{eq:Ep}--\eqref{eq:Eqp} vanishes for all four functionals
simultaneously, leaving
\begin{equation}
E\to2\sum_{i}^{N_e/2}h_{ii}+\sum_{i,j}^{N_e/2}\Big[2J_{ij}-(K_{ij}-K^G_{ij})-(L_{ij}-L^G_{ij})\Big],
\end{equation}
which is exactly the no-pair Dirac--Hartree--Fock energy, Eq.~(66) of
the source paper. This confirms that PNOF5, PNOF7, PNOF7s, and GNOF
all collapse to the same DHF limit, as stated in Sec.~3.6 of the
source paper.

\section{Summary table}
\label{sec:summary}

\begin{center}
\begin{tabular}{@{}lll@{}}
\toprule
Functional & $\PiMat^{\rm inter}_{ij;pq}$ & Explicit $\Dtwo^{j\bar j}_{i\bar i}=\Dtwo^{\bar ii}_{\bar jj}$ \\
\midrule
PNOF5 & $0$ & $0$ \\
PNOF7 & $-\sqrt{n_in_j\hi\hj}$ & $-\tfrac12\sqrt{n_in_j\hi\hj}$ \\
PNOF7s & $-4n_in_j\hi\hj$ & $-2n_in_j\hi\hj$ \\
GNOF & piecewise, Eq.~\eqref{eq:Pi-gnof} & $\tfrac12\PiMat^{\rm inter,\,GNOF}_{ij;pq}$ \\
\bottomrule
\end{tabular}
\end{center}
For every functional, $\Dtwo^{i\bar i}_{\bar jj}=\Dtwo^{j\bar j}_{\bar ii}
=\Dtwo^{\bar ii}_{j\bar j}=\Dtwo^{\bar jj}_{i\bar i}
=-\Dtwo^{j\bar j}_{i\bar i}$, and the ordinary
Coulomb/exchange/spin-flip elements
\eqref{eq:intra-coulomb}--\eqref{eq:diag-common} and
\eqref{eq:inter-coulomb}--\eqref{eq:inter-lflip} are common to all
four.

\section{Conclusion}
\label{sec:conclusion}

We have written down, fully explicitly and separately for PNOF5,
PNOF7, PNOF7s, and GNOF, every reconstructed 2-RDM element
$\Dtwo^{kl}_{ij}$ entering the relativistic no-pair natural-orbital
energy expression. Substituting these explicit elements into Eq.~(64)
and evaluating each of its six blocks in full detail reproduces,
term by term and without any additional assumption, the published
energy expressions Eqs.~(72)--(76) of Sec.~3.5.3 of SciPost Chem.\
\textbf{1}, 004 (2022), and their common no-pair Dirac--Hartree--Fock
limit, Eq.~(66).

\end{document}
